\par\noindent\begin{minipage}{\linewidth}
% !TeX root = ../教材题库_全部条目_转写工作稿.tex
% Source: XJ-B1-C02-PR-P047-Y1177
% 第2章 / 2.2 从函数观点看一元二次方程 / 教材 p.47 / PDF 第54页
% 已对照原页转写题干及公式；待教师复核；未制作答案
% 原题及相邻原解截图（核对用）：../题目原图/XJ-B1-C02-PR-P047-Y1177-01.webp
% 本文件只在转写完成后作为题干使用；OCR 条目尚不是可靠题干。
\begin{bankitem}{p.47 练习 第1题}{第2章 / 2.2 从函数观点看一元二次方程 / 教材 p.47 / PDF 第54页}{XJ-B1-C02-PR-P047-Y1177}{已对照原页转写题干及公式；待教师复核；未制作答案}
利用函数图象判断下列方程有没有实数根，有几个根。
\subq{1}$x^2-4x-3=0$；
\subq{2}$2x(x-2)+3=0$；
\subq{3}$-2x^2+x+3=0$；
\subq{4}$5x^2+2x=3x^2+5$。
\end{bankitem}
\end{minipage}\par\medskip
